% =====================================================================================
%  Chapter 2 — Background and Related Work
%  Thesis: Comparison of Control Strategies for Interaction-Force Aware Multirotor Teams
%
%  DRAFT, 23 August 2026 (writing-track item W.5).
%  REVISED 2026-10-15 (round 2: the hybrid is in scope; one controller for Strategies 0-3).
%  REVISED 2026-10-14 (thesis-revision run, Phase 4). Markers: % NEW, % CHANGED, % DECISION,
%  % VERIFY. Plan rows P2-xx refer to docs/thesis_revision/stage3_revision_plan.md; every changed
%  prior-work sentence was re-checked against the held PDF (records R01-R18 in
%  docs/thesis_revision/stage1_sources.md).
%
%  CITATION DISCIPLINE — read before editing.
%  THIS FILE IS THE SOURCE OF TRUTH for every factual claim about prior work. Verify a claim
%  against the paper itself. Do NOT verify it against, or import it from, any summary, matrix,
%  ledger or review brief in docs/ — those are stale working notes, archived on 2026-08-27 to
%  docs/_archive/citation_notes/, and where they disagree with this file, THIS FILE WINS.
%  Three misattributions in this project came from treating such a note as a source.
%
%  A number belongs here only with what it is a number OF: which axis, which training set,
%  which platform, which experiment, and whether it comes from the abstract or the body.
%  That distinction — not arithmetic — is what every correction pass has turned on.
%
%  Bibliography: docs/references.bib (generated from the arXiv API).
%
%  The Lee et al. citation is RESOLVED: CDC 2010 is the primary reference; its equations were
%  verified against the companion arXiv paper. The remaining \todo{} in §2.9 is a PROVENANCE
%  NOTE recording which audits support the gap claim, not outstanding work.
%
%  Citation baseline: 2026-08-27 (see docs/thesis/CITATION_BASELINE.md).
%  After that date: only NEW cites or EDITED prior-work sentences need re-checking.
% =====================================================================================

\chapter{Background and Related Work}
\label{ch:related}

Multirotors flying in close proximity disturb one another aerodynamically. The disturbance is
large enough to destabilise a vehicle, and conventional practice avoids it by keeping the vehicles
far apart: a Crazyflie~2.0 measures \SI{9}{\centi\metre} rotor-to-rotor, and when the interaction
between vehicles is left unmodelled, planners have adopted conservative separations, Shi et
al.~\cite{shi2020neuralswarm} cite \SI{60}{\centi\metre} from earlier planning work as one such
example. Removing that margin is what makes dense
formation flight possible, and doing so requires a controller that accounts for the interaction
rather than avoiding it.

This chapter surveys how that has been attempted. It is organised as an argument rather than a
catalogue. Section~\ref{sec:rw-physics} establishes what the disturbance is.
Sections~\ref{sec:rw-reactive} and~\ref{sec:rw-predictive} cover online measurement and offline
prediction. Section~\ref{sec:rw-inject} treats where a prediction enters the control law;
Section~\ref{sec:rw-rl} treats learning a corrective action without an explicit force estimate;
Section~\ref{sec:rw-hybrid} treats combining measurement and prediction.
Section~\ref{sec:rw-eval} examines how these methods are evaluated. The field's reporting
conventions make the individual results impossible to compose into a comparison.

Families this thesis does not implement --- direct action learning (Section~\ref{sec:rw-rl}) and
residuals inside a predictive horizon (Section~\ref{sec:rw-inject}) --- are described briefly.
Chapter~\ref{ch:control} states where this thesis's compared strategies end.

% -------------------------------------------------------------------------------------
\section{Multirotor dynamics and baseline control}
\label{sec:rw-baseline}

Each vehicle $i$ in a team is modelled as a rigid body,
%
\begin{align}
  \dot{p}_i &= v_i, &
  m\,\dot{v}_i &= -m g e_3 + f_i R_i e_3 + f_{\mathrm{res},i}, \label{eq:rw-trans}\\
  \dot{R}_i &= R_i \hat{\omega}_i, &
  J\dot{\omega}_i &= -\omega_i \times J\omega_i + \tau_i + \tau_{\mathrm{res},i},
  \label{eq:rw-rot}
\end{align}
%
in a world frame with $e_3$ pointing up, where $f_i$ and $\tau_i$ are the collective thrust and
body torque produced by the rotors. The
terms $f_{\mathrm{res},i}$ and $\tau_{\mathrm{res},i}$ collect every force and moment the nominal
model omits. This chapter is concerned with the case in which they are dominated by the
aerodynamic interaction with neighbouring vehicles. The formulation is developed in
Chapter~\ref{ch:model}.

% NEW (P2-24): background for the reader on rotations and the attitude error.
\paragraph{Rotations and the attitude error.}
The attitude $R_i$ is a rotation matrix: an element of the special orthogonal group $SO(3)$, the
$3\times3$ matrices with $R^\top R = I$ and $\det R = 1$, which maps vectors from the body frame
to the world frame. Its third column $R_i e_3$ is therefore the direction of the thrust in the
world, which is why the thrust appears as $f_i R_i e_3$ in~\eqref{eq:rw-trans}. Position and
attitude together form the pose, an element of $SE(3)$. The hat map $\hat{\omega}$ in~\eqref{eq:rw-rot}
turns a vector into the skew-symmetric matrix with $\hat{\omega} x = \omega \times x$; its inverse
is the vee map $(\cdot)^\vee$. Working with $R$ directly, rather than with three Euler angles,
avoids the singular configurations at which a three-angle description loses a degree of
freedom~\cite{lee2010geometricse3}. A geometric controller measures how far $R$ is from a desired
attitude $R_d$ with the error function $\Psi(R, R_d) = \tfrac12 \operatorname{tr}(I - R_d^\top R)$
and the attitude error vector $e_R = \tfrac12 (R_d^\top R - R^\top R_d)^\vee$~\cite{lee2010geometricse3}.
For two attitudes that differ by a rotation of angle $\theta$ about some axis, $\Psi = 1 - \cos\theta$:
zero when they coincide, one at \SI{90}{\degree} and two at \SI{180}{\degree}. Chapter~\ref{ch:control}
uses both.

The baseline against which interaction-aware methods are measured is a geometric tracking
controller on $SE(3)$, which tracks a position trajectory and a single heading direction
coordinate-free, avoiding the representation singularities of a local attitude
parameterisation~\cite{lee2010geometrictracking}. Its guarantees are regional rather than global:
exponential stability for initial attitude errors below \SI{90}{\degree}, and almost-global
exponential attractiveness below \SI{180}{\degree}. It is the controller three of the
strategies compared in this thesis build upon, and the one used uncompensated as a reference
condition. The attitude error function is that work's, verified equation-by-equation against its
companion arXiv paper, of which version~4 is held~\cite{lee2010geometricse3}; the torque law implemented here is a
\emph{reduced} form of it, and Chapter~\ref{ch:control} states exactly which terms differ.
% VERIFY (CDC 2010 not held): the 90/180 degree regions are checked against the companion v4
%   (Propositions 3 and 4 there); the numbering in the CDC paper is not checked.

A second property of the quadrotor is used throughout: differential flatness. Position and yaw,
together with their derivatives, determine the full state and the inputs, which allows a
feedforward term to be computed directly from a planned trajectory. Tal and
Karaman~\cite{tal2018indi} exploit this to track position and yaw derivatives up to fourth order,
using flatness-derived feedforward for angular rate and angular acceleration in order to follow
jerk and snap.
% NEW (P2-25): what each derivative fixes, from the dynamics above.
The chain follows from~\eqref{eq:rw-trans} and~\eqref{eq:rw-rot} with the residual set to zero. The
desired acceleration fixes the thrust vector $f R e_3 = m(\ddot p + g e_3)$, hence the collective
thrust and the body $z$-axis; the yaw angle then fixes the remaining rotation about that axis, so
that $R$ is determined. Differentiating once more, the jerk fixes the angular velocity, and the
snap fixes the angular acceleration and with it, through~\eqref{eq:rw-rot}, the
torque~\cite{tal2018indi,hsieh2026flatness}. Position and yaw are therefore called the \emph{flat
outputs}. Flatness recurs in Section~\ref{sec:rw-inject} as a structure that a learned
correction can either preserve or destroy.

% -------------------------------------------------------------------------------------
\section{Aerodynamic interaction in multirotor teams}
\label{sec:rw-physics}

% CHANGED (P2-01): the properties carry their sources.
The dominant interaction in a vertically stacked formation is the downwash of the upper vehicle
acting on the lower one. Two properties make it difficult to handle. It is strongly nonlinear,
growing sharply as the vehicles close and decaying quickly with lateral offset; and it is
asymmetric, since the vehicle above is affected far less than the one below. On a rig with two
mounted Crazyflies, the lower vehicle's thrust drops to \SI{65}{\percent} of its hover value at
the closest vertical separation they report, four centre-to-motor distances (about
\SI{13}{\centi\metre}), while the upper one retains
\SI{90}{\percent}; laterally the lower vehicle's thrust is unaltered from an offset of about three
centre-to-motor distances, whereas vertically the effect persists to about nineteen~\cite{kiran2025downwash}.
Measurements on a load stand likewise find a downwash that stays concentrated beneath the vehicle
that produces it, with negligible influence on neighbours outside it apart from small, chaotic
lateral forces~\cite{gielis2023aggregate}.
% CHANGED (closing pass, diff audit): 13 cm marked as derived (Delta z / l = 4, l = 32.5 mm, R11);
%   Gielis' caveat on lateral forces added (R03 Sec. 5).

% CHANGED (P2-02): PIV plane; Gielis sufferer mass; Gielis evaluates its own models.
Recent experimental work characterises the effect directly rather than inferring it from tracking
error. Kiran et al.~\cite{kiran2025downwash} measure forces \emph{and} torques with load cells on
mounted Crazyflies whose roll and yaw are locked and whose motors are driven at a constant hover
voltage, and use particle image velocimetry in a light sheet through the centre of the front two
rotors to resolve the wake, for
a single quadrotor and for a pair; the dynamic cases traverse one vehicle on a rail rather than
flying it. Gielis et al.~\cite{gielis2023aggregate} measure
six-degree-of-freedom exogenous forces on a \SI{700}{\gram} multirotor held on a purpose-built
load stand at hover
thrust, while up to three further vehicles fly freely around it, and report strong nonlinearities
in how those interactions aggregate. The vehicle experiencing the disturbance is therefore
constrained rather than in free flight. Both measure the disturbance as a transducer load rather
than inferring it from a flight controller's residual, actuation is still present in each, at
hover thrust, which makes them, \emph{in our reading}, the natural ground truth against which a
% CHANGED (P2-02): Gielis evaluates its own learned models (R03 Sec. 4.4).
learned model should be checked. Neither paper makes that claim. Gielis et al.\ evaluate their own
learned aggregation models against the stand; neither paper evaluates a model learned from
flight-controller residuals.

Simplified physical models exist. Bauersfeld et al.~\cite{bauersfeld2024airflow} analyse roughly
\SI{16}{\hour} of hover data across six platforms spanning \SIrange{0.23}{6.3}{\kilo\gram}, and
show that once the individual propeller wakes have merged the mean flow is well approximated by a
turbulent jet --- beyond about \num{2.5} motor-to-motor distances $l$ below the vehicle, a length
scale the abstract renders as ``drone-diameters''. That gives a computationally light alternative
to computational fluid dynamics. What it estimates is the \emph{time-averaged magnitude} of the
induced flow \emph{in hover}, not an instantaneous six-degree-of-freedom wrench on a follower. For
the near field the authors report no closed-form first-principles model and point to computational
fluid dynamics instead. Formation flight operates in exactly that near field --- our observation,
not theirs --- which is the practical argument for learning the residual from data rather than
deriving it.
% NEW (P2-27): the model's own closed-loop demonstration.
The model is nonetheless enough for a first-principles predictive compensator in the far field:
used as a feedforward correction of the throttle while one vehicle passes \SI{1}{\metre} and
\SI{2}{\metre} below another, it lowered the height RMSE by a factor of about three to
five~\cite{bauersfeld2024airflow}.

% -------------------------------------------------------------------------------------
\section{Reactive compensation: measuring the residual online}
\label{sec:rw-reactive}

The first family estimates the residual from its effect and requires no prior model of the
interaction. Incremental nonlinear dynamic inversion (INDI) compares measured acceleration with the
acceleration the nominal model predicts for the actuation actually applied, and attributes the
difference to unmodelled forces, which it then cancels incrementally.% NEW (P2-06): terminology.
\footnote{``Reactive'' is used in this thesis for methods that \emph{measure} the residual from its
effect, whether by INDI or by another online estimator such as an $\mathcal{L}_1$-adaptive
module~\cite{hsieh2025onlineadaptation}, an adaptive scheme that estimates the disturbance from the
mismatch between a predicted and a measured state. Part of the model predictive control (MPC)
literature uses the word differently, for any
controller that compensates only the disturbance at the current time step, including feedforward
of a learned model~\cite{chee2024knodedwmpc}; in this thesis such controllers are predictive.}

% NEW (P2-04): how the increment forms the residual, and why the filters must match.
The mechanism can be stated in one line. If the actuation currently applied produces, according to
the nominal model, an acceleration $a_{\mathrm{model}}$, and the accelerometer reports
$a_{\mathrm{meas}}$, the difference is whatever the model omits. INDI does not compute the
command from the model alone but increments the current one: the new command is the current
actuation plus the actuation needed to move the measured acceleration to the desired one, so that
the unmodelled part is carried along inside the measurement rather than being
modelled~\cite{tal2018indi,smeur2017cascadedindi}. The measured acceleration is noisy and must be
low-pass filtered, which delays it. The estimate of the current actuation must then be passed
through an identical filter, so that both terms incur the same phase
lag~\cite{tal2018indi,smeur2017cascadedindi}; otherwise their difference contains a term produced by
the unequal delay rather than by a force.
% CHANGED (closing pass, diff audit): the sources state the requirement; the explanation of the
%   artefact is the thesis's reasoning and is no longer attributed to them. This requirement recurs in Chapter~\ref{ch:control}.

% CHANGED (P2-03): the 4 Hz position updates belong to the wind-tunnel experiment.
Smeur et al.~\cite{smeur2017cascadedindi} cascade an inner angular-acceleration INDI loop with an
outer linear-acceleration INDI loop on a single Parrot Bebop running Paparazzi at
\SI{512}{\hertz}; the position and velocity loops feeding the acceleration reference remain
proportional-derivative rather than a third incremental stage. They demonstrate disturbance
rejection indoors by flying in and out of a \SI{10}{\metre\per\second} open-jet exhaust every
\SI{14}{\second}, and in that same experiment the motion-capture position reached the vehicle at
only \SI{4}{\hertz}. A separate outdoor experiment, taking off and holding position on a consumer
GNSS module in roughly \SI{5.1}{\metre\per\second} average wind, likewise shows that the method
does not depend on frequent, precise position updates. Tal and Karaman~\cite{tal2018indi} use INDI for robust tracking of
linear and angular accelerations in aggressive flight on a single custom \SI{609}{\gram} vehicle,
and state the property that defines the family: no parametric model of the aerodynamic effects is
required. That claim is narrower than it sounds. Mass, inertia, the motor time constant, control
effectiveness and the throttle map are still identified. What is dropped is the drag model.

The same work also exposes the family's cost. Tracking snap requires direct control of body torque,
which those authors achieve using closed-loop motor speed control based on optical encoders on
their own vehicle's motors~\cite{tal2018indi}; encoders are what \emph{that} snap-and-torque path
required, not a general prerequisite for INDI. Reconstructing the applied wrench well enough to
difference it against measurement does, however, generally require some form of rotor-speed
feedback, Smeur et al.\ likewise use motor-speed
measurement~\cite{smeur2017cascadedindi}, which is additional hardware or firmware support that
must be fitted, calibrated and maintained on every vehicle.
% NEW (P2-05): direct evidence for the sensing requirement on this platform class.
On a Crazyflie~2.1, an INDI controller that reconstructed the thrust from the motor commands
instead of measured rotor speeds tracked worse than the uncompensated geometric controller it was
meant to improve~\cite{cobobriesewitz2026lindi}.
A second limitation is structural rather than
practical: the estimate is delayed by the sensing and filtering chain, and it exists only once the
disturbance has already perturbed the vehicle. A reactive controller cannot, even in principle, act
before the disturbance arrives.
% NEW (P2-05): the same limitation as stated by the predictive literature, with its caveat.
Gielis et al.\ make the same point as a motivation for prediction, that reactive controllers cannot
reject the interaction without first observing it~\cite{gielis2023aggregate}. Shankar et al.\ note
that a vehicle docking beneath a leader meets a downwash that a disturbance observer struggles to
estimate because its convergence time is finite~\cite{shankar2023docking}, and Neural-Swarm2 reports
that a simple online estimate-and-compensate scheme quickly crashed the vehicle in their
experiments~\cite{shi2022neuralswarm2}. The last is a single remark without the estimator, its filter
or data, and is evidence about that scheme rather than about INDI.

That the sensing requirement is not fundamental has since been shown. Cobo-Briesewitz et
al.~\cite{cobobriesewitz2026lindi} demonstrate that a neural network can generate smooth
approximations of INDI's residual estimate without rotor-speed measurement, which is discussed in Section~\ref{sec:rw-hybrid}.

% -------------------------------------------------------------------------------------
\section{Predictive compensation: learning the residual from relative state}
\label{sec:rw-predictive}

The second family predicts the residual before it acts, from the relative states of the
neighbours, using a model learned offline from flight data. Because the interaction is a function
of the team's relative configuration, a model with the right inputs can act ahead of the
disturbance rather than after it.

% CHANGED (P2-07): the attitude loop.
Shi et al.~\cite{shi2020neuralswarm} combine a nominal dynamics model with a regularised
permutation-invariant deep neural network that learns multi-vehicle interaction, and use the
learned model as a feed-forward term in a decentralised nonlinear tracking position controller;
the design leaves the attitude controller open, and the experiments use the nonlinear attitude
tracking controller of Lee et al. The learned
force itself is the vertical component \(f_{a,z}\), evaluated onboard the STM32. Training
separate networks on data from two, three and four vehicles respectively, they report worst-case
height error reduced by a factor of two or better, and up to four times smaller in the most
favourable case. In a two-vehicle swapping task the baseline reaches \SI{9.4}{\centi\metre};
the network trained on \emph{four} vehicles brings this to \SI{2.4}{\centi\metre}, while the
network trained on two vehicles reaches \SI{2.7}{\centi\metre}. The swap experiments use vertical
neighbour spacings of \SIrange{0.20}{0.25}{\metre}, with \SI{25}{\centi\metre} stated explicitly
for the three-vehicle case. Generalisation is directional rather than universal: networks trained
on three or four vehicles improved five-vehicle flight, whereas a network trained only on two
vehicles did \emph{not} generalise, performing worse than the baseline on three- and four-vehicle
tests.

% CHANGED (P2-08, P2-26): where each figure is reported; the upgraded-motor qualifier; SN defined.
Neural-Swarm2~\cite{shi2022neuralswarm2} extends this to heterogeneous teams using heterogeneous
deep sets, whose permutation-invariant architecture accepts any number of neighbours. That is a
structural property rather than an accuracy claim: the networks are trained on at most three
vehicles (at most two neighbours), and the paper describes prediction at sixteen robots as
relatively less accurate. Two design choices are worth noting. Spectral normalisation is applied:
it bounds the Lipschitz constant of a network, the largest factor by which its output can change
relative to a change in its input, by constraining the spectral norm of each weight
matrix~\cite{shi2020neuralswarm}. The abstract describes this as yielding stability and generalisation
guarantees; in the body this
rests on Lipschitz bounds together with empirical generalisation and a bounded-residual assumption
used in the proofs, rather than a formal guarantee on unseen data. The learned residual is used in
interaction-aware motion planning as well as in tracking control. The abstract reports up to a
threefold reduction in worst-case tracking error against a baseline nonlinear tracking controller
with the interaction force set to zero; the introduction and its first figure report a
heterogeneous team of sixteen robots flown at a minimum vertical distance of
\SI{24}{\centi\metre}. We do not assume the two describe the same experiment: on the sixteen-robot
flight itself the body reports the maximum \(z\)-error of one small vehicle with upgraded motors,
in one of the vertical rings, falling from \SI{15}{\centi\metre} to \SI{10}{\centi\metre}, a factor
of about \num{1.5}.

% CHANGED (P2-09): where the 56 % is stated; position and velocity; the leader's position.
Subsequent work has focused on making these models cheaper to obtain rather than larger. Smith et
al.~\cite{smith2023so2equivariant} exploit the rotational symmetry of the problem with an
$SO(2)$-equivariant model, and report that five minutes of real-world flight data attains a lower
validation error than learning baselines trained on more than fifteen --- a data-efficiency result
on prediction accuracy rather than a closed-loop tracking one. Deployed online within an optimal
feedback (LQR) controller on two identical vehicles, the model reduces three-dimensional tracking
error by \SI{36}{\percent} and, as the abstract states, vertical tracking error by
\SI{56}{\percent} on average.\footnote{The paper is internally inconsistent on whether the
\SI{36}{\percent} figure is three-dimensional or lateral: we follow the abstract and the
translation-result body, which give three-dimensional and vertical. Its contributions list instead
calls the \SI{36}{\percent} figure lateral. The \SI{56}{\percent} figure appears in the abstract
and the contributions only; the body reports \SI{55}{\percent} for the translation experiment and
about \SI{51}{\percent} for the lemniscate.} On a lateral transect flown beneath the hovering
vehicle they report mean three-dimensional position error falling from \SI{0.154}{\metre} to
\SI{0.098}{\metre}, close to \SI{36}{\percent}, with vertical improvements of \SI{55}{\percent} in
position and \SI{49}{\percent} in velocity. The
platform is not a Crazyflie but a custom quadrotor of roughly \SI{0.26}{\metre} span and
\SI{0.7}{\kilo\gram}, tracked by OptiTrack and controlled from a Raspberry~Pi at \SI{50}{\hertz}
with a \SI{45}{\hertz} network update. The closest \emph{training} stage is about
\SI{0.5}{\metre} --- some two body-lengths --- while the evaluation manoeuvres are flown at
% CHANGED (P2-09): the hovering vehicle is the leader at NED z = -2.5 m (R04 Sec. V-C).
\SIrange{0.6}{0.8}{\metre}, together with a lemniscate flown beneath the leader hovering at a
height of \SI{2.5}{\metre} (\(z = \SI{-2.5}{\metre}\) in their downward-pointing frame). The authors
attribute the efficiency to exploiting the problem's latent geometry rather than learning it from
data.

\label{sec:rw-superposition}
\paragraph{The superposition assumption}

The deep-sets architectures above aggregate per-neighbour terms by summation, which commits them
to a form of superposition. The literature is explicit that this is an approximation.
Shi et al.~\cite{shi2020neuralswarm} observe that with more than two vehicles the aerodynamic
effect is not a simple superposition of each pair. Gielis et al.~\cite{gielis2023aggregate} make
the same point as their motivation, an accurate model of the downwash from one vehicle is
unlikely to suffice for predicting the aggregate from several. Their abstract frames this as
modelling the formation as a graph; the implementation learns the aggregation function with a
sum-pooling deep-sets predictor, contrasted against a linear-summation baseline.
% CHANGED (P2-10): "explore", and K = 3 neighbours is four vehicles in total.
The abstract announces an exploration of how the models generalise when a fourth vehicle is added;
the largest reported case is three neighbours around the stand, four vehicles in total, and the
conclusion cautions that the deep-sets model may be incapable of generalising beyond three
neighbours, so we do not treat generalisation to larger formations as a demonstrated result.

Their measurements refine the picture rather than simply refuting summation. They report that
interactions fall into three regimes depending on the formation and its location in the state
space: a chaotic, noise-dominated regime; a linear regime \emph{well estimated by linear
summation}; and a nonlinear regime requiring a nonlinear model. Superposition is therefore not
uniformly wrong but regime-dependent, which makes the practical question one of how much accuracy
is lost, and where. This is their own load-stand result, obtained independently of Shi et al.'s
free-flight observation cited above.
% NEW (P2-13 context): the same paper is cited in support of summation elsewhere.
Both readings are in use: Hsieh et al.\ cite the same measurements in support of approximating the
aggregate force on a vehicle as a pairwise sum~\cite{hsieh2025onlineadaptation}.

Summed architectures nevertheless remain in wide use, because a controller does not require an
exact model, only one whose remaining residual is smaller than the one it removes. What is not
reported is the size of the penalty that the simplification carries for a deployed model, which is
the question this thesis takes up in the multi-robot results chapter.

% -------------------------------------------------------------------------------------
\section{Injecting the prediction: the controller matters}
\label{sec:rw-inject}

A predicted residual is inert until it enters a control law, and the manner of entry differs by
controller family. Treating ``predictive compensation'' as a single method would obscure
differences that are as large as those between the families themselves.

% CHANGED (P2-11): the closest published instance of a learned residual fed forward into an SE(3) controller.
\paragraph{Feedforward into a geometric controller.}
The simplest injection adds the predicted force to the desired force of the baseline controller,
leaving its structure unchanged. This is the form used with the permutation-invariant models of
Section~\ref{sec:rw-predictive}, whose baseline is a nonlinear tracking
controller~\cite{shi2020neuralswarm,shi2022neuralswarm2}. The closest published instance with a
geometric controller on $SE(3)$ is a configuration in the hardware study of Hsieh et
al.~\cite{hsieh2026flatness}: on two Crazyflie~2.1 vehicles, feeding a learned residual forward into
an $SE(3)$ controller reduced the tracking RMSE by \SI{43.7}{\percent} against the nominal $SE(3)$
controller, with the controller computed off-board.

\paragraph{Feedback linearisation with a flatness-preserving residual.}
A learned correction added to a nominal model can destroy differential flatness, and with it the
applicability of flatness-based planning and control. Yang et al.~\cite{yang2025flatnessresiduals}
show that residuals with a lower-triangular structure preserve both the local flatness of the system
and its original flat output for systems whose nominal model admits a pure-feedback form, and give
a constructive procedure for recovering the diffeomorphism of the augmented system. The residual
there is a \emph{structured correction inside the state derivative} rather than a free
six-degree-of-freedom wrench --- in their worked example it appears only in the velocity equation.
Their validation is in simulation on a planar, two-dimensional quadrotor; there is no hardware, no
three-dimensional vehicle and no multi-robot downwash experiment, and they list a physical quadrotor
as future work. Hsieh et
al.~\cite{hsieh2026flatness} apply the idea to tight formation flight, learning a physics-informed
residual that keeps the joint multi-quadrotor system differentially flat and using the preserved
structure to design a feedback-linearising controller that cancels the disturbance by feedforward.
Their abstract reports a \SI{31}{\percent} reduction in average tracking error against nominal
baselines; the body does not restate that aggregate, reporting instead, on hardware, \SI{24}{\percent}
RMSE and \SI{30}{\percent} maximum-\(z\) improvement for reduced-order-model feedback linearisation (ROM-FBL) over nominal feedback linearisation (FBL), and a further
\SI{29}{\percent} RMSE and \SI{43}{\percent} maximum-\(z\) improvement for learned FBL over
ROM-FBL.

% CHANGED (P2-12): the introduction's wording against the body's.
On hardware they match the approximately \SI{5}{\centi\metre} RMSE and \SI{10}{\centi\metre}
maximum-\(z\) error \emph{reported} by Chee et al.~\cite{chee2024knodedwmpc}, rather than
comparing the two methods on one rig, although their introduction describes the comparison as
made in an identical hardware setup; the head-to-head comparison with nonlinear model predictive
control (NMPC) is in their simulation study, where the flatness-based controller's
maximum-\(z\) error can be worse than that of NMPC.

The \num{20}-fold reduction in solve time (about \SI{1}{\milli\second} against
\SI{27}{\milli\second} at equal residual) is a \emph{simulation} figure. The abstract and the
introduction state an order of magnitude less computation. Training uses \SI{28}{\second} of hardware data
and \SI{18}{\second} in simulation.

Two qualifications matter for our purposes: the hardware
experiments use two Crazyflie~2.1 quadrotors with thrust-upgrade bundles --- very nearly the
platform used in this thesis --- and the controller runs \emph{off-board}, on an Intel~i5 laptop
streaming commands over a Crazyradio link at \SI{400}{\hertz}, rather than on the vehicle. Their
\SI{5}{\milli\second} loop is therefore an abstract-level, off-board figure, and the claim it
supports is that the method is cheap enough to be worth porting, not that it has been shown to run
onboard.

On hardware they also compensate residual \emph{forces} only, commanding thrust and angular
velocity, and delegating residual torques to the vehicle's own rate controller, the same scoping
decision taken in Chapter~\ref{ch:model}. Their simulation study does still model residual torques.

% CHANGED (P2-13): Hsieh et al. 2025 described as what it is; the onboard qualifier.
\paragraph{Inside a model predictive control horizon.}
The residual can instead form part of the prediction model that an optimiser reasons over, so that
the controller plans against a disturbance it expects to meet rather than cancelling the one
present now. Chee et al.~\cite{chee2024knodedwmpc} embed a learned component inside the prediction
constraint of a nonlinear MPC (NMPC); Li et al.~\cite{li2023ndpnmpc} inject a predicted downwash-force
sequence into the optimiser's prediction model. Hsieh et al.~\cite{hsieh2025onlineadaptation} add to
the MPC of Chee et al.\ an $\mathcal{L}_1$-adaptive module that estimates the residual acceleration
online, and fly three Crazyflies in two stacked formations with no compensation, the online
estimate only, the learned prediction only and both. The combination performs best, and for the
middle vehicle the online estimate alone outperforms the learned prediction alone; the authors
also note that the estimate without a downwash model responds more slowly when a vehicle enters or
leaves the wake. Chee et al.\ and Hsieh et al.\ fly Crazyflies but run the
optimisation off-board. Li et al.\ run their controller on a TX2 NX
computer carried by the vehicle, which is not of the Crazyflie class. None of the three runs its
optimiser onboard a Crazyflie-class vehicle. This injection point is not among
the strategies evaluated here --- its onboard feasibility on a Crazyflie-class platform is itself
open --- but it is retained in this survey because its evaluation conditions form part of the
argument in Section~\ref{sec:rw-eval}.

\paragraph{Inside an optimal feedback controller for a docking task.}
A further instance is worth separating because its task differs in kind. Shankar et
al.~\cite{shankar2023docking} embed a learned lumped downwash force online in an infinite-horizon
LQR acting on a feedforward-linearised follower, beneath an inner attitude loop, and use it for
vertical docking: a leader--follower manoeuvre in which the vehicles close to within one or two
body-lengths, so that the disturbance is not incidental to the task but constitutive of it. The
platforms are two custom quadrotors of roughly \SI{0.27}{\metre} diagonal and \SI{0.7}{\kilo\gram},
not the Crazyflie class, and the free-flight dock is performed at a vertical offset of about
\SI{0.36}{\metre} --- close to one body-length --- with a curriculum descending to around
\SI{0.5}{\metre}. The model itself is the $SO(2)$-equivariant one of their prior work, trained
offline from about five minutes of data and queried online on an onboard Raspberry~Pi, with pose
from motion capture; the evaluation is hardware only. Reported accuracy depends on the manoeuvre:
the abstract gives terminal error below \SI{0.06}{\metre}, the body reports terminal vertical error
typically below \SI{0.1}{\metre} for a hovering leader, and for a moving leader below
\SI{0.05}{\metre} with the model against more than \SI{0.2}{\metre} without it. Physical docks
succeeded in eight of ten attempts against a hovering leader at \SI{0.36}{\metre}; against a leader
moving at \SI{0.25}{\metre\per\second} success was limited, which the authors attribute to the
follower's oscillation preventing the gripper from closing, and to the gripper's actuation
latency.
% CHANGED (P2-14): failure attribution as stated in the paper.
It is the closest published setting to the tightest
scenarios of Chapter~\ref{ch:setup}, and the regime in which --- on our reading ---
compensation ceases to be an accuracy improvement and becomes a precondition for the manoeuvre.

% -------------------------------------------------------------------------------------
\section{Learning the compensation directly}
\label{sec:rw-rl}

A third stance declines to model the force at all and learns the corrective action instead. Zhang
et al.~\cite{zhang2024proxfly} train a residual reinforcement-learning module on top of a cascaded
controller, compensating the disturbance and thrust loss caused by downwash without forming an
explicit force estimate at any point. Two properties distinguish the stance sharply from
Section~\ref{sec:rw-predictive}. The policy observes only tracking error, the baseline controller's
command and its own previous action --- no neighbour state and no inter-vehicle communication ---
which removes a bandwidth requirement and a failure mode, at the cost of being unable to anticipate
a neighbour it cannot perceive. And giving up the explicit force estimate removes the
predicted-against-measured comparison that every other learned method here is evaluated by, which
makes failures correspondingly harder to attribute.

This family is not among the strategies evaluated here, for that second reason as much as for
scope: the evaluation built in Chapter~\ref{ch:setup} rests on the residual being a measurable
quantity, and a method that never forms one is assessed on tracking performance alone. It remains
the most natural direction in which to extend the comparison.

% -------------------------------------------------------------------------------------
\section{Hybrid reactive and predictive compensation}
\label{sec:rw-hybrid}

% CHANGED (P2-16): the multi-vehicle MPC instance.
Measurement and prediction are complementary in principle. Prediction is available before the
disturbance acts but is valid only over the configurations its training data covered; measurement
is general but delayed. Combining them is therefore natural. It has been done with INDI for a single
vehicle, described below, and, inside an off-board model predictive controller with an
$\mathcal{L}_1$ estimate, for a team of three (Section~\ref{sec:rw-inject})~\cite{hsieh2025onlineadaptation}.

Cobo-Briesewitz et al.~\cite{cobobriesewitz2026lindi} propose two methods. In the first, LINDI, a
neural network replaces INDI's residual estimate entirely, trained on smoothed data obtained by
spline fitting. In the second, NA-INDI, the total residual is written \(f_a = f_{a,\mathrm{NN}} +
f_{a,\mathrm{INDI}}\): the network predicts a component and the incremental term reasons about the
remaining mismatch. Both are evaluated on a single quadrotor, with and without
a slung payload, on figure-of-eight and circular trajectories, neither of which was in the
training set.

The headline result concerns sensing rather than the hybrid: a network can generate smooth
approximations of INDI's output without requiring rotor-speed sensors at LINDI's runtime, which
removes the hardware prerequisite identified in Section~\ref{sec:rw-reactive} from that method's
deployment --- though rotor speeds are still needed to produce the training labels, and NA-INDI
retains an incremental branch that uses them. Without a payload, residual compensation of any kind
reduces tracking error by around \SI{40}{\percent} relative to the geometric baseline; with a
payload the reductions are smaller, in the range of \SIrange{19}{24}{\percent}. The hybrid's own
advantage is real but modest: without a payload it gives the best performance, though the authors
describe the margin as small, and with a payload the authors describe it as remaining
\emph{similar} to INDI alone. They also observe a slight drop in performance on the circular
trajectory --- for LINDI, in the no-payload case --- attributing it to that trajectory being out of
distribution relative to the training data, a direct instance of the limitation that defines the
predictive family.

% CHANGED (P2-15, P2-16): residual sources as the paper states them; the team instance.
Two features of this setting bear on whether the result transfers. The residual there comes from
the vehicle itself --- the authors list blade flapping, drag and ground forces among the sources of
such residuals --- and, in the payload experiments, from a slung payload, and it is predicted from
the vehicle's own state and motor commands. In a multi-robot team the residual is produced by other
vehicles and is a function of a measurable relative state, which gives the predictive component an
input the single-vehicle setting does not provide. Whether that turns a small margin into a
decisive one has been examined on a team only inside an off-board MPC with an $\mathcal{L}_1$
estimate~\cite{hsieh2025onlineadaptation}, not for an INDI-based hybrid. This thesis flies such a
hybrid on a team (Chapter~\ref{ch:control}).
% CHANGED (round 2, author decision 2).

% -------------------------------------------------------------------------------------
\section{How these methods are evaluated}
\label{sec:rw-eval}

% CHANGED (P2-28): where the Smith 56 % comes from.
The preceding sections describe a field with strong results. Shi et al.~\cite{shi2020neuralswarm}
report up to a fourfold reduction in worst-case height error. Neural-Swarm2~\cite{shi2022neuralswarm2}
reports up to a threefold reduction in worst-case tracking error against a zero-interaction
baseline in one experiment; the sixteen-robot flight is a separate report.
Smith et al.~\cite{smith2023so2equivariant} report \SI{36}{\percent} three-dimensional and, in the
abstract, \SI{56}{\percent} vertical tracking improvements. Hsieh et al.~\cite{hsieh2026flatness} state
\SI{31}{\percent} against nominal baselines in the abstract; the hardware body gives other figures.
Cobo-Briesewitz et al.~\cite{cobobriesewitz2026lindi} report roughly \SI{40}{\percent} for
residual compensation without a payload. Taken individually, each result is credible.

They cannot, however, be composed. Table~\ref{tab:rw-conditions} lists the conditions under which
they were obtained.

% CHANGED (P2-23): caption claim checked against the rows; rows completed.
\begin{table}[t]
  \centering
  \caption{Evaluation conditions across the interaction-force literature. Each result is reported
    against that work's own baseline; rows that share a platform and team size still differ in
    baseline controller, separation or reported metric.}
  \label{tab:rw-conditions}
  \small
  \begin{tabular}{@{}>{\raggedright\arraybackslash}p{0.26\linewidth}>{\raggedright\arraybackslash}p{0.15\linewidth}>{\raggedright\arraybackslash}p{0.22\linewidth}>{\raggedright\arraybackslash}p{0.25\linewidth}@{}}
    \toprule
    Work & Team size & Separation & Reported quantity \\
    \midrule
    Tal and Karaman~\cite{tal2018indi}          & 1                  & n/a (no formation) & aggressive tracking \\
    Smeur et al.~\cite{smeur2017cascadedindi}   & 1                  & n/a & jet / GNSS position error \\
    Shi et al.~\cite{shi2020neuralswarm}        & 2--4 train, 2--5 test & \SIrange{0.20}{0.25}{\metre} vert. & max.\ $|z|$ error \\
    Neural-Swarm2~\cite{shi2022neuralswarm2}    & 16, heterogeneous  & \SI{24}{\centi\metre} vert. & worst-case tracking error\textsuperscript{\dag} \\
    % CHANGED (P2-23): Smith, Gielis, Hsieh 2026, Hsieh 2025 rows completed (R04, R03, R05, R10).
    Smith et al.~\cite{smith2023so2equivariant} & 2, custom \SI{0.7}{\kilo\gram} & \SIrange{0.6}{0.8}{\metre} eval.\ (train.\ to $\approx$\SI{0.5}{\metre}) & 3D and vertical tracking; onboard Pi, mocap \\
    Gielis et al.~\cite{gielis2023aggregate}    & up to 4 total (rig sufferer $+\leq3$ flyers) & test rig, \SI{700}{\gram} sufferer not in free flight & force-prediction accuracy \\
    Yang et al.~\cite{yang2025flatnessresiduals} & 1, planar (sim.\ only) & n/a & simulation only \\
    Hsieh et al.~\cite{hsieh2026flatness}       & 2, Crazyflie~2.1   & \SI{0.3}{\metre} and \SI{0.4}{\metre} vert.\textsuperscript{\ddag} & RMSE and max.\ $z$ error vs.\ nominal FBL/$SE(3)$; off-board \\
    Chee et al.~\cite{chee2024knodedwmpc}       & 2, Crazyflie~2.1   & commanded \SIrange{0.2}{0.4}{\metre}; tight case \(\approx\)\SI{0.08}{\metre} & RMSE / max.\ $z$ vs. nominal MPC; off-board \\
    Li et al.~\cite{li2023ndpnmpc}               & 2, custom          & neighbour-relative & NMPC with predicted force sequence; onboard TX2 NX + mocap \\
    Hsieh et al.~\cite{hsieh2025onlineadaptation} & 3, Crazyflie~2.1 & V-stack and I-stack; \SI{0.2}{\metre} and \SIrange{0.3}{0.4}{\metre} vert. & RMSE / max.\ $z$; estimate, prediction and both, in MPC; off-board \\
    Shankar et al.~\cite{shankar2023docking}    & 2, custom (\(\approx\)\SI{0.27}{\metre}) & dock \(\approx\)\SI{0.36}{\metre} & terminal vert.\ error / dock success; onboard Pi, mocap \\
    Cobo-Briesewitz et al.~\cite{cobobriesewitz2026lindi} & 1 (+ payload) & n/a                 & position tracking error \\
    Zhang et al.~\cite{zhang2024proxfly}        & 2, heterogeneous   & \SI{0.50}{\metre} vert.\ hover/circle; dock to \SI{0.10}{\metre}, release at \SI{0.05}{\metre} & position and attitude error \\
    \bottomrule
  \end{tabular}

  \vspace{0.4em}
  \begin{minipage}{\linewidth}\footnotesize
    \textsuperscript{\dag}~The sixteen-robot flight and the threefold reduction are separate
    reports; on the sixteen-robot flight the body gives a factor of about \num{1.5}.
    \textsuperscript{\ddag}~Hardware tracking comparisons are at commanded
    \(\Delta z = \SI{0.3}{\metre}\) and \SI{0.4}{\metre}; the \SIrange{0.2}{0.5}{\metre} range
    belongs to the simulation and prediction evaluation.
  \end{minipage}
\end{table}

A \SI{36}{\percent} improvement measured on two vehicles in close proximity and a
\SI{31}{\percent} improvement measured in a tight formation are not the same quantity: they differ
in what was flown, on what airframe, at what separation, against what baseline, and under what
metric. Neither can be subtracted from the other, and their ordering carries no information about
which method would prevail if both were flown under the same conditions. The difficulty compounds
where the baselines themselves differ, since an improvement relative to a weaker baseline is
easier to obtain.

% NEW (P2-17): evaluation practice, described as practice rather than fault.
The practice behind the rows has common features. Most results are within-paper ablations: one
controller with the learned term switched on and off, so that each improvement is relative to that
paper's own nominal
controller~\cite{shi2020neuralswarm,shi2022neuralswarm2,chee2024knodedwmpc,li2023ndpnmpc,hsieh2026flatness}.
Metrics converge on a root-mean-square error and a maximum vertical error, but over windows that
differ --- a whole flight, a crossing, a terminal approach. Repetitions are few, between three and
twenty per condition~\cite{li2023ndpnmpc,shi2020neuralswarm,hsieh2026flatness,shankar2023docking,cobobriesewitz2026lindi}.
Separations are mostly reported as commanded; Chee et al.\ report the achieved separation as well,
which in their tightest case differed markedly from the commanded one~\cite{chee2024knodedwmpc}.
Data are split for validation in different ways: by contiguous pieces of flight in
Neural-Swarm2~\cite{shi2022neuralswarm2}, by shuffled samples in Li et al.~\cite{li2023ndpnmpc}, by
held-out trajectories or separations
elsewhere~\cite{cobobriesewitz2026lindi,hsieh2026flatness}. Simulation, a processor in the loop with
a simulator, and flight are different kinds of evidence~\cite{gu2025comparative,yang2025flatnessresiduals},
as are onboard and off-board execution when computational cost is claimed. And the baselines are
rarely tuned for themselves; where tuning for fairness is explicit, it is the exception~\cite{gu2025comparative}.

This is not a criticism of any individual work. Each is internally valid, and there is no
convention in the field that any of them violated. It is a limitation of the literature taken as a
whole: the results are not designed to compose, and no amount of careful reading assembles them
into a ranking. A ranking requires an experiment built for that purpose. That is what distinguishes
a comparison from a survey.

\label{sec:rw-existing-comparisons}
\paragraph{Comparisons that do exist}

% DECISION D-03: the next sentence ("two recent ones bound what remains open") and the paragraph
%   that closes this subsection describe only Gu et al. and Abro et al. Hsieh et al. (2025) flies an
%   online estimate, a learned prediction and both on one three-vehicle rig (Section rw-inject);
%   Zhang et al. compare a learned residual policy with a model-based relative-position compensator
%   on one pair. Options and recommended wording in author_decisions.md. Text left unchanged.
The field does conduct systematic comparisons, and two recent ones bound what remains open.

Gu et al.~\cite{gu2025comparative} benchmark four online-learning uncertainty estimators ---
extreme learning machines, radial-basis-function networks, echo state networks and Gaussian
processes --- inside a single modular framework in which a nominal controller is paired with a
swappable residual-learning module. They treat the uncertainty as a lumped additive wrench
\((f_a, \tau_a)\) recovered by inverting the discrete Newton--Euler dynamics --- analogous to the
residual acceleration used here, though not identical to it a priori. They evaluate estimation accuracy, tracking
performance and computational cost, and reach the kind of conclusion a comparison should: not a
winner, but a set of trade-offs, with Gaussian processes most accurate yet limited by cubic
training complexity, the lightweight parametric models fast but sensitive to initialisation, and
echo state networks the most balanced.

Two things separate it from the present work. Its subject is a \emph{single} quadrotor under
generic uncertainty, and all four methods learn \emph{online}; it therefore compares estimators
\emph{within} one family rather than the families of Sections~\ref{sec:rw-reactive}
to~\ref{sec:rw-hybrid}. And its evaluation is in simulation throughout: the reported hardware
component is the learning modules executing on an embedded computer in the loop with a simulator,
which measures latency rather than flight, and the authors list validation on a real quadrotor as
future work.

% CHANGED (P2-21): the "no PDF" comment that stood here concerned Abro et al.; this paragraph is
%   about Gu et al. and was checked against the held PDF (R18 Sec. 1, 5.1).
Its approach to fairness is nonetheless directly relevant. They tune the parametric models --- the
extreme learning machine, the radial-basis-function network and the echo state network --- to their
own optima by Bayesian optimisation under otherwise consistent conditions, so that each operates
near its best achievable performance. This thesis answers the same objection --- that an observed
difference might reflect tuning effort rather than method --- differently: its reactive,
predictive and hybrid strategies are switch settings of one controller with one shared gain set,
frozen before the comparison and the training flights. Section~\ref{sec:ctrl-fairness} states what
that does and does not guarantee.
% CHANGED (round 2, author decision 6): per-strategy tuning replaced by one shared gain set.

% UNVERIFIED (review P-010): no PDF of Abro et al. in docs/papers; check this paragraph against the paper.
% VERIFY (P2-22): every clause of the next paragraph (methods compared, close-proximity and downwash
%   motivation, number of vehicles, simulation or hardware) against the full text.
Abro et al.~\cite{abro2025fractional} compare fractional-order PID, fractional-order LQR and
integral state feedback on close-proximity flight, explicitly motivated by downwash and by the
excessive safety distances that neglecting it imposes.\footnote{Cited for the methods compared and
for that motivation; the abstract's quantitative improvement figure and its claim about larger
swarms are not used here.} This is a controller comparison in our setting: multiple vehicles, and
downwash-aware. Of the three, only integral state feedback is classical; the fractional-order laws
are not. All three are nonetheless feedback regulators: they do not present a learned or
first-principles predictor of the interaction force in the sense of the works above, so the
reactive, predictive and hybrid families remain uncompared against one another.

% CHANGED (round 2): "per-strategy gains frozen after tuning" -> "gains frozen" in the requirements below.
% DECISION D-03: the paragraph below states the gap as "the families ... have not been placed under
%   common conditions"; Hsieh et al. (2025) places all three information sources under common
%   conditions inside one MPC. Text left unchanged.
Together these delimit the gap precisely. It is not that close-proximity control is uncompared, nor
that the field lacks the methodology. It is that the specific families defined by \emph{how they
obtain information about the interaction force}, measure it, predict it, or both, have not
been placed under common conditions.

Such an experiment has requirements that follow directly from the failure modes above: a single
airframe; identical trajectories and formations across all methods; gains frozen before data
collection, so that a difference cannot be a tuning artefact; a single definition
of the residual, measured identically for every controller
including those that do not use it; and a null condition in which no interaction is present, to
establish the measurement floor and so distinguish a real effect from noise. Chapter~\ref{ch:setup}
describes the resulting design.

% -------------------------------------------------------------------------------------
\section{Summary and gap}
\label{sec:rw-gap}

A controller can obtain information about the interaction force in two ways. It can measure the
residual online, which requires no prior model and generalises to disturbances never seen, but is
delayed and, in its established form, depends on rotor-speed sensing. Or it can predict the
residual from the relative states of its neighbours using a model learned offline, which acts ahead
of the disturbance but is valid only where the training data reached. A third approach learns the
corrective action directly, forgoing an explicit force estimate and, in at least one instance,
inter-vehicle communication as well. Representative methods exist for each, and each is
individually validated.

Three gaps follow from this survey.

% DECISION D-03: first gap. Recommended replacement (Stage 2 Sec. 3): "To the best of our knowledge,
%   no study has compared rotor-speed-based reactive compensation (INDI) with learned per-neighbour
%   feedforward prediction in onboard controllers on one airframe, across a fixed library of
%   interaction regimes, with the residual logged under every controller and the transfer of the
%   summed model from two to three vehicles quantified. The closest study compares an online
%   L1 estimate, a learned downwash model and their combination inside one off-board MPC on three
%   vehicles in two stacked formations~\cite{hsieh2025onlineadaptation}." Text left unchanged.
First, and principally, to the best of our knowledge the reactive, predictive and hybrid families
have not been evaluated against one another under identical conditions. Systematic comparisons in
adjacent territory exist --- of online-learning estimators on a single vehicle, and of classical
regulators for close-proximity swarms (Section~\ref{sec:rw-existing-comparisons}) --- but none
spans the families distinguished by how they obtain the interaction force. The individual results are strong; what the field lacks is not strong results
but comparable ones, and its reporting conventions --- improvements against each work's own
baseline, on its own platform, at its own separations --- make the comparison impossible to
assemble after the fact. This thesis compares the uncompensated baseline, reactive compensation
(INDI), predictive compensation (geometric+NN) and their hybrid (INDI+NN), and FBL+NN if it is
implemented.
% CHANGED (round 2, author decision 2): the hybrid is in scope; the gap statement itself is D-03.

% CHANGED (P2-20): minimal factual correction; restructuring of this gap is part of D-03.
Second, the hybrid of measurement and prediction has been demonstrated with INDI only for a single
vehicle, where the residual is self-induced and its advantage was described as small without a
payload and as leaving performance similar to the reactive method with
one~\cite{cobobriesewitz2026lindi}. On a team it has been flown inside an off-board model predictive
controller with an $\mathcal{L}_1$ estimate, where the combination performed
best~\cite{hsieh2025onlineadaptation}. An INDI-based hybrid on a team, in which the residual is
larger and is a function of a measurable relative state, remains untested; this thesis evaluates
one on a team.
% CHANGED (round 2, author decision 2).

% CHANGED (D-03 context, factual): the summed model has been deployed on three vehicles.
Third, the summation used to aggregate per-neighbour predictions is known to be an
approximation~\cite{shi2020neuralswarm,gielis2023aggregate}, yet the magnitude of the penalty it
carries for a deployed model, evaluated on teams larger than those it was trained on, is not
reported. A summed model trained on two-vehicle data has been flown on three
vehicles~\cite{hsieh2025onlineadaptation}, but the penalty was not isolated.

% DECISION D-03: the provenance note below records audits that did not classify Hsieh et al.
%   (2025) as spanning the families; the author should revise it together with the gap. Unchanged.
\todo[color=green!25]{\textbf{Provenance, not pending work.} Supported by two audits: an arXiv keyword sweep, and a
cross-database sweep over OpenAlex, Crossref and DBLP including a citation-graph scan of 157 works
citing the Neural-Swarm line. Both found adjacent comparisons but none spanning these families.
Neither is exhaustive; keep the ``to the best of our knowledge'' phrasing.}

Chapter~\ref{ch:model} formalises the residual wrench these methods estimate or predict, and
Chapter~\ref{ch:control} describes the strategies compared in this thesis and the point at
which each injects its knowledge of the interaction into the control law.
